
I mainly follow the discussion in \cite{Carlip:1998qg}. For more
detailed discussion and having the same notation one can refer to
\cite{carrollSpacetimeGeometryIntroduction2014}.

\subsection{Vielbein and Spin Connections}

We can a GR theory using a first order formalism, which use frame
(aka vielbein) and spin connection as basic variables instead of
metric. This is especially useful in 2+1 dimensional gravity, since
it can be reformulated as a Chern-Simons theory.
We define the frame fields
\defi{
  Frame Fields are a set of orthogonal vectors in the tangent space $
  e^{\mu}{}_a $ or the inverse one form $ e_{\mu}^a $, which satisfy:
  \begin{align}
    g^{\mu\nu}e_\mu{}^ae_\nu{}^b=\eta^{ab},\quad
    e_\mu{}^ae_\nu{}^b\eta_{ab}=g_{\mu\nu}.
  \end{align}
  Here $ e^\mu{}_a $ is the inverse of $ e_\mu{}^a $, i.e.,
  \begin{align}
    {e^\mu}_a e_\mu{}^b = \delta_a^b, \quad e_\mu{}^a {e^\nu}_a =
    \delta_\mu^\nu.
  \end{align}
}
We can show that $ e^a = e_\mu{}^a d x^\mu $ forms a orthonormal
basis of the cotangent space and $ e_a =  {e^\mu}_a \partial_\mu $
forms a orthonormal basis of the tangent space. The coeficients can
also be understood as transformation matrices between coordinate
basis and orthonormal basis.

This set of vectors definately forms a basis of the tangent space,
thus we can consider the connections in this basis. For example, the
Christoffel symbols can be written as:
\begin{align}
  \nabla_\mu \partial_{\nu} = \Gamma_{\mu\nu}^\rho \partial_\rho,
  \quad \text{or}\quad \nabla_\mu dx^{\nu} = -\Gamma_{\mu\rho}^\nu dx^\rho
\end{align}
Thus, a analogue is to define a connection $ \omega_\mu{}^a{}_b $ as
\defi{
  Spin Connection $ \omega_\mu{}^a{}_b $ is defined as:
  \begin{align}
    \nabla_\mu{e_\nu}^a = -\omega_\mu{}^a{}_b{e_\nu}^b
  \end{align}
}
\rmk{
  The covariant derivative here $ \nabla_\mu e_\nu{}^a $ means that
  we only view $   e_\nu{}^a$ as a dual vector in the cotangent
  space, with componet $ \mu $ of the coordinate basis. Sometimes you
  may encounter the notation:
  \begin{align}
    \nabla_\mu{e_\nu}^a=\partial_\mu{e_\nu}^a-\Gamma_{\mu\nu}^\lambda{e_\lambda}^a+\omega_\mu{}^a{}_b
    {e_\nu}^b
  \end{align}
  This is when we view it as a (1,1) tensor in a mixed basis, with
  the coordinate basis in the cotangent space and the orthonormal
  basis in the tangent space.
}

The covariant derivative of a vector is a rank (0,2) tensor, thus we
know that $ \omega_\mu{}^a{}_b $ transforms as a one form in the $
\mu $ index. Thus:
\begin{itemize}
  \item Both $ e^a = e_\mu{}^a dx^\mu $ and $ \omega^a{}_b =
    \omega_\mu{}^a{}_b dx^\mu $ can be understood as one forms. The
    inverse of the frame fields can be written as a vector field $
    e_a = {e^\mu}_a \partial_\mu $.
  \item We define the lowering and raising of the spin connection
    indices $ \omega_{ab} = \eta_{ac}\omega^c{}_b $.
\end{itemize}
We can prove that if the theory is metric compatible, then the spin
connection 1-form is antisymmetric in $ a,b $ indices:
\begin{align}
  \omega_{ ab} = -\omega_{ ba}.
\end{align}

\subsection{Cartan Structure Equations}

Using the above definitions, we hope to express important geometric
quantities such as Riemann curvature Tensor. Moreover, we also notice
that the frame fields and the spin connections are not independent,
just like the Christoffel symbol can be represented by the metric. We want to:
\begin{itemize}
  \item Represent the spin connection in terms of the frame fields or
    at least find a relation between them.
  \item Represent the Riemann curvature Tensor in terms of the spin connection.
\end{itemize}
These two goals can be achieved by the Cartan structure equations:

\thm{
  Cartan's first structure equation

  If the theory is torsion free, then the frame fields and the spin
  connection satisfy:
  \begin{align}
    D_{\omega}e^a \equiv de^a+\omega^a{}_b\wedge e^b=0,
  \end{align}
  $ D_{\omega}e^a $ is called the gauge covariant derivative of $ e^a
  $ with respect to the spin connection $ \omega $.
}
This can be proved directly from the definition of the covariant
derivative and the definition of the spin connection. The torsion
free condition makes the lower indices of the Christoffel symbol
symmetric, thus it cancels out when doing a exterior derivative.

\subsection{Local Lorentz Transformations}

\YL{[This is important see carrol]}
